\section*{Chapitre \ref{ch:s1:intraday-machine-learned-alphas} --- Alphas intrajournaliers appris par apprentissage automatique}

\begin{solution}{exo:s1:intraday-machine-learned-alphas:1}
La corrélation vaut $\sqrt{0.25} = 0.5$, donc une prévision bien calibrée a un écart type de $0.5 \times 0.69 = 0.34$ ticks~: la plupart des prévisions valent un tiers de
tick ou moins, bien en dessous d'une fourchette d'un tick.
\end{solution}

\begin{solution}{exo:s1:intraday-machine-learned-alphas:2}
Les deux cibles partagent 29 de leurs 30 secondes, donc un modèle testé sur la seconde $t + 1$ après un entraînement sur la seconde $t$ a vu presque toute la cible de test~: le
test mesure la mémoire, non la prévision. L'entraînement et le test doivent être séparés d'au moins l'horizon (purge et embargo), ici par des sessions entières.
\end{solution}

\begin{solution}{exo:s1:intraday-machine-learned-alphas:3}
Une fourchette~: une prévision d'au moins un tick dans le sens de la transaction, avant frais.
\end{solution}

\begin{solution}{exo:s1:intraday-machine-learned-alphas:4}
Un ordre passif gagne la fourchette lorsqu'il est exécuté~; la prévision choisit le côté vers lequel le prix a le plus de chances d'aller, donc les exécutions sont moins soumises à la sélection adverse
que des exécutions aléatoires. Le modèle d'exécution est optimiste~: il exécute l'ordre dès que le marché traverse son prix, comme s'il était en tête de
file.
\end{solution}

\begin{solution}{exo:s1:intraday-machine-learned-alphas:5}
Il a ajusté du bruit~: avec peu de signal à deux minutes, les arbres ont trouvé dans les sessions d'entraînement des motifs qui ne se reproduisent pas, et la prévision ajoute
de la variance sans pouvoir prédictif, de sorte que le R-carré hors échantillon est inférieur à zéro.
\end{solution}

\begin{solution}{exo:s1:intraday-machine-learned-alphas:6}
Le même flux d'ordres déplace davantage le prix lorsque le carnet d'ordres est mince~: diviser par la profondeur rend la variable comparable entre les moments et les instruments, et plus proche
de la stationnarité, comme le suggère la relation linéaire de Cont, Kukanov et Stoikov, dont la pente est inversement proportionnelle à la profondeur.
\end{solution}

\begin{solution}{exo:s1:intraday-machine-learned-alphas:7}
Sans les variations passées du prix milieu, le R-carré à cinq secondes vaut 18.6~\% pour la régression ridge (contre 21.1~\%) et 24.4~\% pour les arbres (contre 24.6~\%)~: presque rien n'est perdu.
Les variables de flux d'ordres et de déséquilibre contiennent déjà la tendance à court terme, qui sur le tape synthétique est pilotée par le flux lui-même.
\end{solution}

\begin{solution}{exo:s1:intraday-machine-learned-alphas:8}
La justesse directionnelle ne dit rien de la taille des mouvements par rapport à la fourchette~: une prévision à une seconde vaut une fraction de tick, et un ordre au marché paie
une demi-fourchette à l'entrée et autant à la sortie. Elle ne doit être négociée que de façon passive ou pas du tout, et la justesse doit être mesurée hors échantillon avec une latence réaliste.
\end{solution}

\begin{solution}{pb:s1:intraday-machine-learned-alphas:1}
\begin{enumerate}
\item Une prévision de quelques secondes à quelques heures, construite à partir des données de marché et exécutée dans la journée~; l'intervalle sur lequel le mouvement est prévu.
\item La variation du prix milieu sur 5, 30 et 120 secondes, à chaque seconde, sur vingt sessions d'une heure~; une fourchette d'environ 1.07 ticks.
\item Déséquilibre de profondeur, fourchette, déséquilibre du flux d'ordres sur 1, 5 et 30 secondes, volume signé sur 5 et 30, rendements propres sur 5 et 30, distance au microprix~;
les variables ne doivent pas changer lorsque le futur est retiré.
\item Les entrées de flux d'ordres battent les états bruts du carnet d'ordres~; l'horizon effectif vaut environ deux variations moyennes de prix.
\item Régression ridge et soixante arbres boostés de profondeur deux.
\item Entraînement sur les sessions 1--14, test sur les sessions 15--20, afin que les cibles qui se recouvrent ne créent aucune fuite d'information.
\item 21.1~\% et 24.6~\% à 5 secondes, 5.1~\% à 30, inférieur à zéro à 120~; IC~: 0.33--0.34, 0.23--0.28, 0.07.
\item Les prix milieux du tape synthétique suivent une tendance de quelques secondes, portée par son flux d'ordres.
\item Choisir le type d'ordre et le prix à partir de la prévision~: ne croiser que lorsque la prévision dépasse la fourchette.
\item Agressif $-0.60$, $-0.48$, $-0.87$~; passif 0.37, 0.33, $-0.02$~; couplé 0.33, 0.22, $-0.24$ ticks par action.
\item La fourchette qu'ils paient dépasse les mouvements qu'ils prévoient.
\item Les exécutions ont lieu dès que le marché traverse le prix de l'ordre~: en tête de file.
\item \textbf{Résultat principal.} R-carré hors échantillon de 24.6~\%, 5.1~\% et $-1.5~\%$ à 5, 30 et 120 secondes (arbres boostés)~; après la fourchette, le trading agressif
de chaque prévision perd 0.60 ticks par action à cinq secondes, tandis que les entrées passives gagnent 0.37 et la règle couplée 0.33.
\item Un modèle unique entraîné sur toutes les actions prédit mieux que des modèles propres à chaque action, y compris hors échantillon, sur des actions qu'il n'a pas vues.
\item Normaliser les variables et les cibles par le tick, la fourchette et la profondeur de chaque instrument, et entraîner un seul modèle sur l'ensemble.
\item Les exécutions passives seraient moins nombreuses et plus sélectionnées défavorablement~; les P\&L passif et couplé diminueraient.
\item La prévision se dégrade en quelques secondes, donc un retard dans les données ou les ordres détruit la part prévisible.
\item Exécution pilotée par l'alpha, qui réduit les coûts sur les transactions que la société réalise de toute façon.
\item Petite~: une fraction de la taille au contact, pour chaque valeur.
\item Ce qu'elle économise ou gagne après la fourchette, dans l'exécution pour laquelle elle est conçue.
\end{enumerate}
\end{solution}

\begin{solution}{iq:s1:intraday-machine-learned-alphas:1}
Déséquilibre au meilleur niveau du carnet et déséquilibre de profondeur, déséquilibre du flux d'ordres sur plusieurs fenêtres, volume signé des transactions, la fourchette, la distance du microprix au prix milieu, les rendements
récents, et les mêmes pour les instruments liés~; tout étant normalisé par la profondeur ou la volatilité.
\end{solution}

\begin{solution}{iq:s1:intraday-machine-learned-alphas:2}
Découper dans le temps avec un écart au moins égal à l'horizon (purge), suivi d'un embargo, ou par sessions entières~; ne jamais mélanger aléatoirement, et présenter des résultats en walk-forward.
\end{solution}

\begin{solution}{iq:s1:intraday-machine-learned-alphas:3}
Non~: l'utiliser pour décider où placer des ordres passifs et quand s'en abstenir, ce qui rapporte ou économise une partie de la fourchette~; elle ne devient inutile que si même les exécutions passives
sont sélectionnées défavorablement au-delà de ce que la prévision permet de récupérer.
\end{solution}

\begin{solution}{iq:s1:intraday-machine-learned-alphas:4}
Aplatir les arbres en tableaux de seuils et de valeurs de feuilles, mettre à jour les variables de façon incrémentale à chaque message, garder tout en cache, et évaluer
dans du code compilé (C++ ou Rust) sans allocation~; ou pré-calculer des tables de correspondance pour les quantiles des variables.
\end{solution}

\begin{solution}{iq:s1:intraday-machine-learned-alphas:5}
Parce qu'elle est négociée des milliers de fois par jour sur de nombreux instruments~: un petit edge par transaction, s'il survit aux coûts, s'additionne~; et les rendements à court horizon sont
surtout du bruit, si bien qu'un R-carré de 1~\% peut correspondre à une corrélation de 0.1, un coefficient d'information (IC) élevé.
\end{solution}

\begin{solution}{iq:s1:intraday-machine-learned-alphas:6}
Étant donné $f$, $E[y \mid f] = \rho(\sigma_y/\sigma_f)f$. Négocier $\operatorname{sign}(f)$ lorsque $|f| > c$ rapporte $E[\rho(\sigma_y/\sigma_f)|f| - c \mid |f| > c]$ par
transaction~; avec $f$ gaussien, $E[|f| \mid |f| > c] = \sigma_f\,\phi(c/\sigma_f)/(1 - \Phi(c/\sigma_f))$, donc le P\&L croît linéairement avec
$\rho$ et n'est positif que lorsque $\rho\sigma_y$ multiplié par cette moyenne de queue dépasse $c$~: une corrélation faible exige un seuil élevé, et les transactions sont rares.
\end{solution}
